We present in turn two representative calculi for
the two dominant styles of graded coeffects:
Linear Base (Section~\ref{sec:linear-base}) and
Graded Base (Section~\ref{sec:graded-base}).

\subsection{Linear Base}
\label{sec:linear-base}

As discussed in the introduction, linear-base systems take
linear $\lambda$-calculus as their core. Such systems include
the core quantitative coeffect calculus of \citet{brunel2014core},
 the combined graded system for effects and coeffects of \citet{DBLP:conf/icfp/GaboardiKOBU16},
$\lambda\mathcal{R}$ of \citet{wood2020linear},
the core of Granule by \citet{DBLP:journals/pacmpl/OrchardLE19},
and associated work~\cite{DBLP:journals/corr/abs-2112-14966}.

The term syntax of our idealised system, Linear Base, is as follows:
%
\begin{align*}
    \SYSTEMnt{t} & ::= \SYSTEMmv{x} \mid  \SYSTEMnt{t_{{\mathrm{1}}}} \,  \SYSTEMnt{t_{{\mathrm{2}}}}  \mid  \lambda  \SYSTEMmv{x}  .  \SYSTEMnt{t}  \mid  \textcolor{coeffectColor}{[}  \SYSTEMnt{t}  \textcolor{coeffectColor}{]}  \mid  \mathsf{let} \, \textcolor{coeffectColor}{[}  \SYSTEMmv{x}  \textcolor{coeffectColor}{]} =  \SYSTEMnt{t_{{\mathrm{1}}}}  \, \mathsf{in} \,  \SYSTEMnt{t_{{\mathrm{2}}}} 
  \tag{terms}
\end{align*}
%
Terms comprise variables, applications, and abstraction as standard. Graded modal introduction
(formally \emph{promotion}, informally \emph{boxing}) is via $ \textcolor{coeffectColor}{[}  \SYSTEMnt{t}  \textcolor{coeffectColor}{]} $, which for now we can
think of as wrapping a term in a constructor. Graded modal elimination eliminates this constructor
via a \emph{let} binding, unwrapping the boxed value, which we refer to as \emph{unboxing}.

\noindent
Linear Base types include a linear function space, a
\emph{graded modality} which can be understood as generalising Linear Logic's exponential
modality $!$, and a base type $ \mathrm{K} $ so that types are well-founded:
%
\begin{align*}
  \SYSTEMnt{A} & ::=  \SYSTEMnt{A}  \multimap  \SYSTEMnt{B}  \mid  \textcolor{coeffectColor}{\square_{ \SYSTEMnt{r} } }  \SYSTEMnt{A}  \mid  \mathrm{K} 
  \tag{types}
\end{align*}
%
Here $\SYSTEMnt{r}$ is drawn from a pre-ordered
semiring $(\mathcal{R}, \cdot, 1, +, 0, \sqsubseteq)$
(where addition $+$ and multiplication $\cdot$ must
be monotonic wrt. to $\sqsubseteq$)
parameterising the calculus.

An intuition for $ \textcolor{coeffectColor}{\square_{ \SYSTEMnt{r} } }  \SYSTEMnt{A} $ is that it describes some capability $r$
with which terms of type $A$ contained inside can be used. For example,
for the natural numbers semiring $\mathbb{N}$ then $ \textcolor{coeffectColor}{\square_{ \SYSTEMsym{2} } }  \SYSTEMnt{A} $ denotes two copies
of an $A$, or in essence that its contents can be used twice.

Typing judgments are of the form $ \Gamma  \vdash_{\textsc{l} }  \SYSTEMnt{t}  :  \SYSTEMnt{A} $, relating a context $\Gamma$, term $\SYSTEMnt{t}$, and type $\SYSTEMnt{A}$. Contexts comprise linear assumptions
$ \SYSTEMmv{x}  :  \SYSTEMnt{A} $ and graded assumptions $ \SYSTEMmv{x}  : \textcolor{coeffectColor}{[}  \SYSTEMnt{A} {\textcolor{coeffectColor}{]_{ \SYSTEMnt{r} } } } $
with grade $\SYSTEMnt{r} \in \mathcal{R}$, of the form:
%
\begin{align*}
  \tag{contexts}
  \Gamma ::=  \emptyset  \mid  \Gamma ,   \SYSTEMmv{x}  :  \SYSTEMnt{A}   \mid  \Gamma ,   \SYSTEMmv{x}  : \textcolor{coeffectColor}{[}  \SYSTEMnt{A} {\textcolor{coeffectColor}{]_{ \SYSTEMnt{r} } } }  
\end{align*}
%
Contexts can be \textit{added} if they are disjoint in their linear assumptions and
using semiring addition to combine grades of any graded assumptions shared
between two contexts:
%
\begin{definition}[Context addition]
    For contexts $\Gamma_{{\mathrm{1}}}, \Gamma_{{\mathrm{2}}}$, then $\Gamma_{{\mathrm{1}}}  \SYSTEMsym{+}  \Gamma_{{\mathrm{2}}}$ is a partial operation
    computing the combined
    context providing \textit{contraction} on graded assumptions. Context addition is specified:
\begin{align*}
\Gamma + \emptyset = \Gamma \quad
 (   \Gamma_{{\mathrm{1}}} ,   \SYSTEMmv{x}  : \textcolor{coeffectColor}{[}  \SYSTEMnt{A} {\textcolor{coeffectColor}{]_{ \SYSTEMnt{r} } } }    )   \SYSTEMsym{+}   (   \Gamma_{{\mathrm{2}}} ,   \SYSTEMmv{x}  : \textcolor{coeffectColor}{[}  \SYSTEMnt{A} {\textcolor{coeffectColor}{]_{ \SYSTEMnt{s} } } }    )  & =   (  \Gamma_{{\mathrm{1}}}  \SYSTEMsym{+}  \Gamma_{{\mathrm{2}}}  )  ,   \SYSTEMmv{x}  : \textcolor{coeffectColor}{[}  \SYSTEMnt{A} {\textcolor{coeffectColor}{]_{  \SYSTEMnt{r}  \SYSTEMsym{+}  \SYSTEMnt{s}  } } }   \\
\emptyset + \Gamma = \Gamma \hspace{5.4em}  (   \Gamma_{{\mathrm{1}}} ,   \SYSTEMmv{x}  : \textcolor{coeffectColor}{[}  \SYSTEMnt{A} {\textcolor{coeffectColor}{]_{ \SYSTEMnt{r} } } }    )   \SYSTEMsym{+}  \Gamma_{{\mathrm{2}}} & =   (  \Gamma_{{\mathrm{1}}}  \SYSTEMsym{+}  \Gamma_{{\mathrm{2}}}  )  ,   \SYSTEMmv{x}  : \textcolor{coeffectColor}{[}  \SYSTEMnt{A} {\textcolor{coeffectColor}{]_{ \SYSTEMnt{r} } } }   & \mathrm{if}\ x \not\in \mathsf{dom}(\Gamma_{{\mathrm{2}}}) \\
 (   \Gamma_{{\mathrm{1}}} ,   \SYSTEMmv{x}  :  \SYSTEMnt{A}    )   \SYSTEMsym{+}  \Gamma_{{\mathrm{2}}} & =   (  \Gamma_{{\mathrm{1}}}  \SYSTEMsym{+}  \Gamma_{{\mathrm{2}}}  )  ,   \SYSTEMmv{x}  :  \SYSTEMnt{A}   & \mathrm{if}\ x \not\in \mathsf{dom}(\Gamma_{{\mathrm{2}}})
\end{align*}
That is, we union the disjoint parts of the contexts, and for the non-disjoint parts
of the contexts we merge with semiring addition. If the two contexts share linear assumptions
or if the two contexts share graded assumptions that do not agree on the type, then
the addition is undefined (in practical terms, this amounts to a linearity error).
Contexts are unordered, i.e., \emph{exchange} is permitted.
\end{definition}
%
We define two predicates for denoting contexts containing only graded assumptions
and those containing graded assumptions all with the same grade:
%
\begin{definition}[Graded contexts]
\label{dfn:graded-contexts}
  A context $\Gamma$ has property $ \mathrm{graded}( \Gamma ) $ if it contains only graded
  assumptions, and $ \mathrm{graded}( \Gamma , \textcolor{coeffectColor}{ \SYSTEMnt{r} }) $ if it contains only graded assumptions which are
  of grade $\SYSTEMnt{r}$. These predicates are defined inductively:
\begin{align*}
  \begin{array}{c}
\SYSTEMdrulePredicateGradednil{}
\quad
\SYSTEMdrulePredicateGradedgrad{}
\qquad
\SYSTEMdrulePredicateGradedAtnil{}
\quad
\SYSTEMdrulePredicateGradedAtgrad{}
\end{array}
\end{align*}
\end{definition}
%
\noindent
Lastly, contexts which contain only graded assumptions can be \emph{scaled} by
a grade:
%
\begin{definition}[Scalar multiplication]
  A context $\Gamma$ where $ \mathrm{graded}( \Gamma ) $ can be \emph{scaled}
  $ \textcolor{coeffectColor}{ \SYSTEMnt{r}  \cdot}  \Gamma $, defined:
  %
  \begin{align*}
         \textcolor{coeffectColor}{ \SYSTEMnt{r}  \cdot}   \emptyset   \triangleq  \emptyset 
    \qquad
         \textcolor{coeffectColor}{ \SYSTEMnt{r}  \cdot}   (   \Gamma' ,   \SYSTEMmv{x}  : \textcolor{coeffectColor}{[}  \SYSTEMnt{A} {\textcolor{coeffectColor}{]_{ \SYSTEMnt{s} } } }    )   \triangleq   (   \textcolor{coeffectColor}{ \SYSTEMnt{r}  \cdot}  \Gamma'   )  ,   \SYSTEMmv{x}  : \textcolor{coeffectColor}{[}  \SYSTEMnt{A} {\textcolor{coeffectColor}{]_{   \SYSTEMnt{r}  \cdot  \SYSTEMnt{s}   } } }  
  \end{align*}
  %
\end{definition}
%
%
\begin{figure}[t]
  \raggedright\noindent\framebox{$ \Gamma  \vdash_{\textsc{l} }  \SYSTEMnt{t}  :  \SYSTEMnt{A} $} \quad Typing rules
\begin{gather*}
\begin{align*}
  \begin{array}{c}
\SYSTEMdruleLinvar{}
\quad
\SYSTEMdruleLinabs{}
\quad
\SYSTEMdruleLinapp{}
\\[1.25em]
\SYSTEMdruleLinweak{}
\quad
\SYSTEMdruleLinder{}
\quad
\SYSTEMdruleLinpr{}
\\[1.25em]
\SYSTEMdruleLinlet{}
\qquad
\SYSTEMdruleLinapprox{}
\end{array}
\end{align*}
\end{gather*}
\caption{Linear Base typing.}
\label{fig:linear-base-types}
\end{figure}

% \newpage
Figure~\ref{fig:linear-base-types} defines typing rules. The \textit{variable}
($\SYSTEMRenameRuleLinvar{}$) and \textit{abstraction} ($\SYSTEMRenameRuleLinabs{}$) rules are standard for the
linear $\lambda$-calculus. The \textit{application} ($\SYSTEMRenameRuleLinapp{}$) rule is close to that of the
$\lambda$-calculus, but the separate contexts of the premises are added in the conclusion.
\textit{Weakening} ($\SYSTEMRenameRuleLinweak{}$) allows unused assumptions to be added, as long as they
are graded at the zero of the semiring. \textit{Dereliction} ($\SYSTEMRenameRuleLinder{}$) treats a linear
assumption as a graded assumption at the one element of the semiring, which connects the linear and
non-linear worlds. \textit{Promotion} ($\SYSTEMRenameRuleLinpr{}$) introduces grading to terms, making
use of scalar context multiplication to propagate the resource requirements to the assumptions. The
\textit{let} or \textit{unboxing} rule ($\SYSTEMRenameRuleLinlet{}$) eliminates an $\SYSTEMnt{r}$-graded
modality, making the unboxed value of $\SYSTEMnt{t_{{\mathrm{1}}}}$ available in the body of $\SYSTEMnt{t_{{\mathrm{2}}}}$ according to the capabilities
afforded by $\SYSTEMnt{r}$. \textit{Approximation} ($\SYSTEMRenameRuleLinapprox{}$) allows treating an $\SYSTEMnt{r}$-graded assumption as $\SYSTEMnt{s}$-graded if $ \SYSTEMnt{r}  \, \textcolor{coeffectColor}{\sqsubseteq} \,  \SYSTEMnt{s} $. Note the direction of the graded
inequality in the premise is according to \citet{petricek2014coeffects}. The opposite polarity is
also used in the literature, for example by \citet{brunel2014core}, without any semantic difference.


\begin{example}
\label{exm:nat-eq}
The natural numbers semiring $(\mathbb{N}, \ast, 1, +, 0, \equiv)$ whose ordering
is equality (i.e., a discrete ordering) captures a notion of exact usage.
For example, the following captures the typing of a linear function $f$ that
takes two inputs of type $\SYSTEMnt{A}$. The inputs are provided by first eliminating
an assumption $ \SYSTEMmv{x}  :   \textcolor{coeffectColor}{\square_{ \SYSTEMsym{2} } }  \SYSTEMnt{A}  $ into a graded assumption $ \SYSTEMmv{y}  : \textcolor{coeffectColor}{[}  \SYSTEMnt{A} {\textcolor{coeffectColor}{]_{ \SYSTEMsym{2} } } } $ which is then
used twice to apply $f$:
%
\begin{align*}
\inferrule*[right=\SYSTEMRenameRuleLinlet{}]
{
\ldots
   \inferrule*[right=\SYSTEMRenameRuleLinapp{}]
    {\inferrule*[right=\SYSTEMRenameRuleLinder{}]
      {\inferrule*[right=\SYSTEMRenameRuleLinapp{}]
        {\ldots}
        {    \SYSTEMmv{f}  :    \SYSTEMnt{A}  \multimap  \SYSTEMnt{A}   \multimap  \SYSTEMnt{B}    ,   \SYSTEMmv{y}  :  \SYSTEMnt{A}    \vdash_{\textsc{l} }   \SYSTEMmv{f} \,  \SYSTEMmv{y}   :   \SYSTEMnt{A}  \multimap  \SYSTEMnt{B}  }}
        {    \SYSTEMmv{f}  :    \SYSTEMnt{A}  \multimap  \SYSTEMnt{A}   \multimap  \SYSTEMnt{B}    ,   \SYSTEMmv{y}  : \textcolor{coeffectColor}{[}  \SYSTEMnt{A} {\textcolor{coeffectColor}{]_{ \SYSTEMsym{1} } } }    \vdash_{\textsc{l} }   \SYSTEMmv{f} \,  \SYSTEMmv{y}   :   \SYSTEMnt{A}  \multimap  \SYSTEMnt{B}  }
      \,
      \inferrule*[right=\SYSTEMRenameRuleLinder{}]
           {\inferrule*[right=\SYSTEMRenameRuleLinvar{}]
             {\quad}
             {   \SYSTEMmv{y}  :  \SYSTEMnt{A}    \vdash_{\textsc{l} }  \SYSTEMmv{y}  :  \SYSTEMnt{A} }
           }
           {   \SYSTEMmv{y}  : \textcolor{coeffectColor}{[}  \SYSTEMnt{A} {\textcolor{coeffectColor}{]_{ \SYSTEMsym{1} } } }    \vdash_{\textsc{l} }  \SYSTEMmv{y}  :  \SYSTEMnt{A} }
     \hspace{-2em}}
     {    \SYSTEMmv{f}  :    \SYSTEMnt{A}  \multimap  \SYSTEMnt{A}   \multimap  \SYSTEMnt{B}    ,   \SYSTEMmv{y}  : \textcolor{coeffectColor}{[}  \SYSTEMnt{A} {\textcolor{coeffectColor}{]_{ \SYSTEMsym{2} } } }    \vdash_{\textsc{l} }    \SYSTEMmv{f} \,  \SYSTEMmv{y}  \,  \SYSTEMmv{y}   :  \SYSTEMnt{B} \hspace{-6em}}
\hspace{-5em}}
{    \SYSTEMmv{f}  :    \SYSTEMnt{A}  \multimap  \SYSTEMnt{A}   \multimap  \SYSTEMnt{B}    ,   \SYSTEMmv{x}  :   \textcolor{coeffectColor}{\square_{ \SYSTEMsym{2} } }  \SYSTEMnt{A}     \vdash_{\textsc{l} }     \mathsf{let} \, \textcolor{coeffectColor}{[}  \SYSTEMmv{y}  \textcolor{coeffectColor}{]} =  \SYSTEMmv{x}  \, \mathsf{in} \,  \SYSTEMmv{f}  \,  \SYSTEMmv{y}  \,  \SYSTEMmv{y}   :  \SYSTEMnt{B} }
\end{align*}
%
\end{example}

\begin{example}[None-one-tons~\cite{McBride2016,quantitative-type-theory}, a.k.a. the Linearity Semiring]
  \label{exm:none-one-tons}
  Linearity can be captured with the semiring $\mathcal{R} = \{\SYSTEMsym{0}, \SYSTEMsym{1},  \omega \}$, with the required laws for $0$ and $1$, and $\SYSTEMsym{1}  \SYSTEMsym{+}  \SYSTEMsym{1} =  \omega $ and $ \omega   \SYSTEMsym{+}  \SYSTEMnt{r} =  \omega  = \SYSTEMnt{r}  \SYSTEMsym{+}   \omega $ and $  \omega   \cdot   \omega   =  \omega $, and ordering $ \SYSTEMsym{0}  \, \textcolor{coeffectColor}{\sqsubseteq} \,   \omega  $ and $ \SYSTEMsym{1}  \, \textcolor{coeffectColor}{\sqsubseteq} \,   \omega  $. This captures linear reasoning in a graded modal setting~\cite{DBLP:journals/pacmpl/AbelB20,DBLP:journals/pacmpl/AbelDE23},
  with $\SYSTEMsym{1}$ denoting linearity and $ \omega $ for
  non-linearity, and $\SYSTEMsym{0}$ providing a slightly finer-grained
  view when non-linearly discarding.
\end{example}

\begin{example}[Security~\cite{DBLP:journals/pacmpl/AbelB20,DBLP:conf/icfp/GaboardiKOBU16,mycroftfest2024}]
\label{exm:security-lattice}
Consider a semi-lattice of security levels $\mathcal{L} = \{\mathsf{Lo},
\mathsf{Hi}\}$ with ordering such that $\mathsf{Hi} \succeq \mathsf{Lo}$.
This forms a semiring with multiplication as join $\cdot = \vee$
with unit $1 = \mathsf{Lo}$, and addition as meet $+ = \wedge$
with unit $0 = \mathsf{Hi}$. We can type a function which maps
low-security to high-security values, using approximation where we set $\sqsubseteq \; = \; \succeq$:

\vspace{-1.5em}
\begin{align*}
\inferrule*[right=\SYSTEMRenameRuleLinapprox{}]
{
 \inferrule*[right=\SYSTEMRenameRuleLinpr{}]
 {
 \inferrule*[right=\SYSTEMRenameRuleLinder{}]
  {
    \inferrule*[right=\SYSTEMRenameRuleLinvar{}]
      {\quad}
      {   \SYSTEMmv{x}  :  \SYSTEMnt{A}    \vdash_{\textsc{l} }  \SYSTEMmv{x}  :  \SYSTEMnt{A} }
  }
  {   \SYSTEMmv{x}  : \textcolor{coeffectColor}{[}  \SYSTEMnt{A} {\textcolor{coeffectColor}{]_{  \mathsf{Lo}  } } }    \vdash_{\textsc{l} }  \SYSTEMmv{x}  :  \SYSTEMnt{A} }
 }
 {   \SYSTEMmv{x}  : \textcolor{coeffectColor}{[}  \SYSTEMnt{A} {\textcolor{coeffectColor}{]_{  \mathsf{Hi}  } } }    \vdash_{\textsc{l} }   \textcolor{coeffectColor}{[}  \SYSTEMmv{x}  \textcolor{coeffectColor}{]}   :   \textcolor{coeffectColor}{\square_{  \mathsf{Hi}  } }  \SYSTEMnt{A}   \qquad \mathsf{Hi} \succeq \mathsf{Lo}}
}
{   \SYSTEMmv{x}  : \textcolor{coeffectColor}{[}  \SYSTEMnt{A} {\textcolor{coeffectColor}{]_{  \mathsf{Lo}  } } }    \vdash_{\textsc{l} }   \textcolor{coeffectColor}{[}  \SYSTEMmv{x}  \textcolor{coeffectColor}{]}   :   \textcolor{coeffectColor}{\square_{  \mathsf{Hi}  } }  \SYSTEMnt{A}  }
\end{align*}
%
A function mapping from $ \mathsf{Hi} $ to $ \mathsf{Lo} $ values is derivable but cannot
use its input. This \emph{non-interference} property is proved by
\citet{DBLP:journals/pacmpl/AbelB20} and \citet{mycroftfest2024}.
\end{example}

\begin{example}[Product]
  \label{exm:product}
Given two pre-ordered semirings $\mathcal{R}$ and $\mathcal{S}$,
their product $\mathcal{R} \times \mathcal{S}$ is a semiring, with
zero element $\langle 0 , 0 \rangle \in (\mathcal{R} \times \mathcal{S})$,
one element $\langle 1 , 1 \rangle \in (\mathcal{R} \times \mathcal{S})$, and
operations and ordering acting pointwise. This is a standard
algebraic construction which has been used in graded systems to combine
analyses (e.g.~\cite{DBLP:journals/pacmpl/OrchardLE19}).

Section~\ref{subsec:linearity-preserving} uses this construction to
compose the linearity semiring (Example~\ref{exm:none-one-tons}) with some other
given semiring, thereby making explicit the linearity analysis inherent to Linear Base.

\citet{DBLP:journals/pacmpl/BianchiniDGZ23} present the \emph{smash product} construction
as a variation on the product construction which has some preconditions and filters
out the elements of the product that have either the left or the right $ = 0$,
instead providing a single canonical $\SYSTEMsym{0}$ element.
\end{example}

\begin{lemma}[Admissibility of substitution]
  There are two admissible substitutions:
  \begin{enumerate}
  \item (Linear) If $ \Gamma_{{\mathrm{1}}}  \vdash_{\textsc{l} }  \SYSTEMnt{t_{{\mathrm{1}}}}  :  \SYSTEMnt{A} $ and $  \Gamma_{{\mathrm{2}}} ,   \SYSTEMmv{x}  :  \SYSTEMnt{A}    \vdash_{\textsc{l} }  \SYSTEMnt{t_{{\mathrm{2}}}}  :  \SYSTEMnt{B} $
    then $ \Gamma_{{\mathrm{1}}}  \SYSTEMsym{+}  \Gamma_{{\mathrm{2}}}  \vdash_{\textsc{l} }   [  \SYSTEMnt{t_{{\mathrm{1}}}}  /  \SYSTEMmv{x}  ]  \SYSTEMnt{t_{{\mathrm{2}}}}   :  \SYSTEMnt{B} $;

  \item (Graded) If $ \Gamma_{{\mathrm{1}}}  \vdash_{\textsc{l} }  \SYSTEMnt{t_{{\mathrm{1}}}}  :  \SYSTEMnt{A} $ and $ \mathrm{graded}( \Gamma_{{\mathrm{1}}} ) $ and
    $  \Gamma_{{\mathrm{2}}} ,   \SYSTEMmv{x}  : \textcolor{coeffectColor}{[}  \SYSTEMnt{A} {\textcolor{coeffectColor}{]_{ \SYSTEMnt{r} } } }    \vdash_{\textsc{l} }  \SYSTEMnt{t_{{\mathrm{2}}}}  :  \SYSTEMnt{B} $
    then $  \textcolor{coeffectColor}{ \SYSTEMnt{r}  \cdot}  \Gamma_{{\mathrm{1}}}   \SYSTEMsym{+}  \Gamma_{{\mathrm{2}}}  \vdash_{\textsc{l} }   [  \SYSTEMnt{t_{{\mathrm{1}}}}  /  \SYSTEMmv{x}  ]  \SYSTEMnt{t_{{\mathrm{2}}}}   :  \SYSTEMnt{B} $.
  \end{enumerate}
  \label{lemma:linear-base-sub-admiss}
\end{lemma}


\begin{remark}
In the literature (e.g.~\citet{brunel2014core,DBLP:conf/icfp/GaboardiKOBU16,DBLP:journals/pacmpl/OrchardLE19}),
the graded context predicates (\ref{dfn:graded-contexts}) are often presented instead as partial identity functions on contexts:
written $[ \Gamma ]$, the function is identity when the context contains only
graded assumptions and is undefined otherwise, preventing a rule from
being applied (and similarly $[ \Gamma ]_r$ for contexts
containing only assumptions graded by $r$).
We chose to make the predicates explicit for clarity.
\end{remark}

\paragraph{Operational Semantics}

Figure~\ref{fig:linear-base-ops-and-eqns} (first group) gives the
operational semantics, which is that of the call-by-name
$\lambda$-calculus with an additional congruence rule
for graded modal elimination and $\beta$-reduction of graded modal
introduction followed by elimination. %Note the rule for
%evaluating inside of a graded modal introduction in keeping with
%the call-by-name style.

Note that the operational semantics we consider for each calculus in this paper represents the standard
models given in the literature in question, with progress and
preservation theorems, but these do not track specific resource
behaviour operationally (see discussion of resource-aware operational
models in Section~\ref{sec:raos}).

Appendix~\refappendix{app:definitions} gives the definition of syntactic
substitution which is standard.

\begin{figure}[t]
\raggedright\noindent\framebox{$ \SYSTEMnt{t}  \rightsquigarrow_{\textsc{l} }  \SYSTEMnt{t} $} \quad Operational semantics (Call-By-Name)
\begin{align*}
  \begin{array}{c}
  \SYSTEMdruleSemLinbeta{}
  \quad
  \SYSTEMdruleSemLincongAppL{}
  \\[1em]
  \SYSTEMdruleSemLinbetaBox{}
  \quad
  \SYSTEMdruleSemLincongLetL{} \\[1em]
  \end{array}
\end{align*}
\bigskip

\raggedright\noindent\framebox{$ \SYSTEMnt{t}  \,\equiv_{\textsc{l} }\,  \SYSTEMnt{t'} $}
\hspace{1em} Equational theory (congruence rules omitted).
\begin{gather*}
\begin{align*}
  \begin{array}{c}
\SYSTEMdruleLinEqbeta{} \;\;
\SYSTEMdruleLinEqeta{} \;\;
\SYSTEMdruleLinEqbetaBox{} \\[1em]
\SYSTEMdruleLinEqetaBox{} \quad
\SYSTEMdruleLinEqletCommBox{}  \\[1em]
\SYSTEMdruleLinEqletCommOne{} \, \SYSTEMdruleLinEqletCommTwo{}
  \end{array}
\end{align*}
\end{gather*}
\caption{Linear Base - Operational semantics and equations}
\label{fig:linear-base-ops-and-eqns}
\end{figure}


\paragraph{Equational theory}

Figure~\ref{fig:linear-base-ops-and-eqns} (second group) defines an equational theory, giving the key rules.
For brevity, we omit the congruence rules which apply to all
sub-term positions. Note that $ \rightsquigarrow_{\textsc{l} }  \subseteq  \equiv_{\textsc{l} } $.

\begin{definition}[Freshness]
  We make use of freshness here and throughout. It is defined as follows:
  \begin{align*}
    x_1, \ldots, x_n \# t \triangleq \{x_1, \ldots, x_n\} \cap \mathsf{fv}(t) = \emptyset
  \end{align*}
\end{definition}

\subsection{Graded Base}
\label{sec:graded-base}
The next calculus, Graded Base, is representative of the coeffect systems of
\citet{DBLP:conf/icalp/PetricekOM13,petricek2014coeffects},
the generalised bounded linear type system of \citet{ghica2014bounded},
the simply-typed core of work by \citet{DBLP:journals/pacmpl/AbelB20}, \citet{grad}, and \citet{DBLP:journals/pacmpl/AbelDE23},
the OO-based approach of \citet{DBLP:journals/pacmpl/BianchiniDGZS22},
and others.
The syntax of terms is that of the $\lambda$-calculus:
%
\begin{align*}
    \SYSTEMnt{t} & ::= \SYSTEMmv{x} \mid  \SYSTEMnt{t_{{\mathrm{1}}}} \,  \SYSTEMnt{t_{{\mathrm{2}}}}  \mid  \lambda  \SYSTEMmv{x}  .  \SYSTEMnt{t} 
  \tag{terms}
\end{align*}
%
We use $t$ to also range over Graded Base terms but disambiguate which
calculus terms belong to via the context of their later use (in translations).
Type comprise graded functions and a base type $ \mathrm{K}  $:
%
\begin{align*}
  \SYSTEMnt{A} & ::=  \SYSTEMnt{A}  \xrightarrow{\textcolor{coeffectColor}{ \SYSTEMnt{r} } }  \SYSTEMnt{B}  \mid  \mathrm{K} 
  \tag{types} \\
  \Delta & ::=  \emptyset  \mid  \Delta ,   \SYSTEMmv{x}  :_{\textcolor{coeffectColor}{ \SYSTEMnt{r} } }  \SYSTEMnt{A}  
  \tag{contexts}
\end{align*}
%
Instead of a linear function type and graded modality, Graded Base has just a \emph{graded function type}, where $r \in \mathcal{R}$. Contexts contain only graded assumptions which we instead denote as $ \SYSTEMmv{x}  :_{\textcolor{coeffectColor}{ \SYSTEMnt{r} } }  \SYSTEMnt{A} $.

Figure~\ref{fig:graded-base} (first group) gives the typing with judgments of the form $ \Delta  \vdash_{\textsc{g} }  \SYSTEMnt{t}  :  \SYSTEMnt{A} $. Note the singleton context graded at $\SYSTEMsym{1}$
in the \textit{variable} rule. The \textit{abstraction} rule introduces its binding into the context
graded at $\SYSTEMnt{r}$, per the graded function arrow. In the \textit{application} rule, the function
arrow propagates the function's resource requirements onto its argument $\SYSTEMnt{t_{{\mathrm{2}}}}$ via scalar
multiplication on $\Delta_{{\mathrm{2}}}$ (i.e., the context of $\SYSTEMnt{t_{{\mathrm{2}}}}$). Weakening scales by $0$ all the grades in an additional context $\Delta'$.
%
\begin{figure}[t]
\raggedright\noindent\framebox{$ \Delta  \vdash_{\textsc{g} }  \SYSTEMnt{t}  :  \SYSTEMnt{A} $} \quad Typing rules
\begin{align*}
\begin{array}{c}
\SYSTEMdruleGradvar{}
\quad
\SYSTEMdruleGradweak{}
\quad
\SYSTEMdruleGradapprox{}
\\[1.25em]
\SYSTEMdruleGradabs{}
\quad
\SYSTEMdruleGradapp{}
\end{array}
\end{align*}
\bigskip

\raggedright\noindent\framebox{$ \SYSTEMnt{t}  \rightsquigarrow_{\textsc{g} }  \SYSTEMnt{t'} $} \quad Operational semantics
\vspace{-1em}
\begin{align*}
  \begin{array}{c}
  \SYSTEMdruleSemGrdbeta{}
  \quad
  \SYSTEMdruleSemGrdcongAppL{}
  \end{array}
\end{align*}
\bigskip

\raggedright\noindent\framebox{$ \SYSTEMnt{t}  \equiv_{\textsc{g} }  \SYSTEMnt{t'} $} \quad Equational theory (congruence rules omitted)
\vspace{-1em}
\begin{align*}
  \begin{array}{c}
\SYSTEMdruleGradEqbeta{} \quad
\SYSTEMdruleGradEqeta{}
  \end{array}
\end{align*}
\caption{Graded Base system}
\label{fig:graded-base}
\end{figure}

Graded substitution is admissible, akin to the second clause
of Linear Base substitution (Lemma~\ref{lemma:linear-base-sub-admiss}):
%
\begin{lemma}[Admissibility of substitution]
   If $ \Delta_{{\mathrm{1}}}  \vdash_{\textsc{g} }  \SYSTEMnt{t_{{\mathrm{1}}}}  :  \SYSTEMnt{A} $ and
    $  \Delta_{{\mathrm{2}}} ,   \SYSTEMmv{x}  :_{\textcolor{coeffectColor}{ \SYSTEMnt{r} } }  \SYSTEMnt{A}    \vdash_{\textsc{g} }  \SYSTEMnt{t_{{\mathrm{2}}}}  :  \SYSTEMnt{B} $ then\\
    $  \textcolor{coeffectColor}{ \SYSTEMnt{r}  \cdot}  \Delta_{{\mathrm{1}}}   \SYSTEMsym{+}  \Delta_{{\mathrm{2}}}  \vdash_{\textsc{g} }   [  \SYSTEMnt{t_{{\mathrm{1}}}}  /  \SYSTEMmv{x}  ]  \SYSTEMnt{t_{{\mathrm{2}}}}   :  \SYSTEMnt{B} $.
\end{lemma}

\begin{example}
The simply-typed $\lambda$-calculus is recovered
by instantiating Graded Base with the single-point semiring $\mathcal{R} = \{\bullet\}$.
\end{example}
%
\begin{example}
The higher-order function $h =  \lambda  \SYSTEMmv{f}  .    \lambda  \SYSTEMmv{x}  .   \SYSTEMmv{f} \,  \SYSTEMsym{(}   \SYSTEMmv{f} \,  \SYSTEMmv{x}   \SYSTEMsym{)}    $ mentioned in the introduction, composing a function with
itself twice, has the following derivation for its body:
%
\begin{gather*}
\inferrule*[right=\SYSTEMRenameRuleGradapp{}]
{\inferrule*[right=\SYSTEMRenameRuleGradvar{}]{\;}{   \SYSTEMmv{f}  :_{\textcolor{coeffectColor}{ \SYSTEMsym{1} } }   (   \SYSTEMnt{A}  \xrightarrow{\textcolor{coeffectColor}{ \SYSTEMnt{r} } }  \SYSTEMnt{A}   )     \vdash_{\textsc{g} }  \SYSTEMmv{f}  :   \SYSTEMnt{A}  \xrightarrow{\textcolor{coeffectColor}{ \SYSTEMnt{r} } }  \SYSTEMnt{A}  }
%
 \inferrule*[right=\SYSTEMRenameRuleGradapp{}]
 {\inferrule*[right=\SYSTEMRenameRuleGradvar{}]{\;}{   \SYSTEMmv{f}  :_{\textcolor{coeffectColor}{ \SYSTEMsym{1} } }   (   \SYSTEMnt{A}  \xrightarrow{\textcolor{coeffectColor}{ \SYSTEMnt{r} } }  \SYSTEMnt{A}   )     \vdash_{\textsc{g} }  \SYSTEMmv{f}  :   \SYSTEMnt{A}  \xrightarrow{\textcolor{coeffectColor}{ \SYSTEMnt{r} } }  \SYSTEMnt{A}  }
 %
  \inferrule*[right=\SYSTEMRenameRuleGradvar{}]{\;}{   \SYSTEMmv{x}  :_{\textcolor{coeffectColor}{ \SYSTEMsym{1} } }  \SYSTEMnt{A}    \vdash_{\textsc{g} }  \SYSTEMmv{x}  :  \SYSTEMnt{A} }}
  {    \SYSTEMmv{f}  :_{\textcolor{coeffectColor}{ \SYSTEMsym{1} } }   (   \SYSTEMnt{A}  \xrightarrow{\textcolor{coeffectColor}{ \SYSTEMnt{r} } }  \SYSTEMnt{A}   )    ,   \SYSTEMmv{x}  :_{\textcolor{coeffectColor}{ \SYSTEMnt{r} } }  \SYSTEMnt{A}    \vdash_{\textsc{g} }   \SYSTEMmv{f} \,  \SYSTEMmv{x}   :  \SYSTEMnt{A}  }
\hspace{-3em}}
{    \SYSTEMmv{f}  :_{\textcolor{coeffectColor}{ \SYSTEMsym{(}  \SYSTEMsym{1}  \SYSTEMsym{+}  \SYSTEMnt{r}  \SYSTEMsym{)} } }   (   \SYSTEMnt{A}  \xrightarrow{\textcolor{coeffectColor}{ \SYSTEMnt{r} } }  \SYSTEMnt{A}   )    ,   \SYSTEMmv{x}  :_{\textcolor{coeffectColor}{ \SYSTEMsym{(}   \SYSTEMnt{r}  \cdot  \SYSTEMnt{r}   \SYSTEMsym{)} } }  \SYSTEMnt{A}    \vdash_{\textsc{g} }   \SYSTEMmv{f} \,  \SYSTEMsym{(}   \SYSTEMmv{f} \,  \SYSTEMmv{x}   \SYSTEMsym{)}   :  \SYSTEMnt{A} }
\end{gather*}
%
\end{example}
All the semirings mentioned in Section~\ref{sec:linear-base} (Linear Base) can also be used with
Graded Base.


\paragraph{Operational Semantics and Equational theory}

Figure~\ref{fig:graded-base} (second group) gives the
operational semantics, which is that of the CBN
$\lambda$-calculus. Appendix~\refappendix{app:definitions} gives the definition of syntactic
substitution which is standard.

Figure~\ref{fig:graded-base} (third group) defines an equational theory
for the calculus, giving the key rules but omitting the congruence
rules which apply to all sub-term positions. Note $ \rightsquigarrow_{\textsc{g} }  \, \subseteq \,  \equiv_{\textsc{g} } $.
